\documentclass[11pt,a4paper]{article}

% Packages
\usepackage[utf8]{inputenc}
\usepackage[T1]{fontenc}
\usepackage{amsmath,amssymb,amsthm}
\usepackage{mathtools}
\usepackage{geometry}
\usepackage{hyperref}
\usepackage{cleveref}
\usepackage{enumitem}
\usepackage{booktabs}
\usepackage{array}
\usepackage{xcolor}

% Page geometry
\geometry{margin=1in}

% Theorem environments
\theoremstyle{plain}
\newtheorem{theorem}{Theorem}[section]
\newtheorem{lemma}[theorem]{Lemma}
\newtheorem{proposition}[theorem]{Proposition}
\newtheorem{corollary}[theorem]{Corollary}

\theoremstyle{definition}
\newtheorem{definition}[theorem]{Definition}
\newtheorem{example}[theorem]{Example}

\theoremstyle{remark}
\newtheorem{remark}[theorem]{Remark}

% Custom commands
\newcommand{\R}{\mathbb{R}}
\newcommand{\Cspace}{\mathcal{C}}
\newcommand{\Pspace}{\mathcal{P}}
\newcommand{\Sclass}{\mathcal{S}}
\newcommand{\Vocab}{\mathcal{V}}
\newcommand{\ip}[2]{\langle #1, #2 \rangle}
\newcommand{\norm}[1]{\| #1 \|}
\newcommand{\proj}[1]{\Pi_{#1}}
\newcommand{\Attn}{\mathrm{Attn}}
\newcommand{\FFN}{\mathrm{FFN}}

% Epistemic markers
\newcommand{\proven}{\textsc{[Proven]}}
\newcommand{\derived}{\textsc{[Derived]}}
\newcommand{\conjectured}{\textsc{[Conjectured]}}
\newcommand{\empirical}{\textsc{[Empirical]}}

% Title
\title{\textbf{Provenance-Gated Architectures}\\[0.5em]
\large Formal Foundations for Injection-Resistant Language Models}

\author{
Aaron Ben-Shalom$^{1,*}$ \and Claude$^{2}$\\[0.5em]
\small $^1$Independent Researcher, Ember Research Lab\\
\small $^2$Large Language Model, Anthropic\\[0.5em]
\small $^*$Corresponding author: abenshalom305@gmail.com
}

\date{January 2026}

\begin{document}

\maketitle

\begin{abstract}
We present a formal mathematical framework for Provenance-Gated Architecture (PGA), a proposed modification to transformer-based language models that provides architectural immunity to prompt injection attacks. By partitioning the residual stream into orthogonal content and provenance subspaces, and constraining layer updates to the content subspace, we prove that input source information is preserved through arbitrary network depth. We establish dimensionality bounds for robust signature encoding under noise, and prove that content-only adversaries cannot forge authorization for privileged actions. These results provide theoretical foundations for a new class of injection-resistant AI systems.

\medskip
\noindent\textbf{Keywords:} prompt injection, formal verification, transformer architecture, orthogonal subspaces, provenance preservation, adversarial robustness, AI security
\end{abstract}

\tableofcontents
\newpage

%==============================================================================
\section{Introduction}
%==============================================================================

Current transformer architectures treat all tokens in the context window as computationally equivalent. This design choice, while enabling flexible in-context learning, creates a fundamental security vulnerability: the model cannot architecturally distinguish between trusted instructions and untrusted input. All distinctions must be made through reasoning over content, which can be manipulated.

This paper establishes formal foundations for an alternative architecture in which input \emph{provenance}---the pathway through which tokens entered the network---is preserved as an invariant property through all layers of computation. We prove that under specified constraints, this provenance information cannot be forged through content manipulation, enabling provably secure action gating.

\subsection{Contributions}

\begin{enumerate}
    \item A formal definition of \emph{content-constrained layers} that preserve orthogonal provenance information
    \item A proof that provenance signatures are invariant through arbitrary network depth
    \item Dimensionality bounds for robust signature encoding under realistic noise conditions
    \item An \emph{impossibility theorem} showing content-only adversaries cannot forge privileged authorization
    \item Tightness results establishing necessary conditions for security guarantees
\end{enumerate}

%==============================================================================
\section{Preliminaries}
%==============================================================================

\subsection{Notation}

Let $d$ denote the model dimension (residual stream width). We use $\ip{\cdot}{\cdot}$ for inner product and $\norm{\cdot}$ for Euclidean norm. For a subspace $S \subseteq \R^d$, let $\proj{S}: \R^d \to S$ denote orthogonal projection onto $S$.

\subsection{Transformer Architecture}

We consider decoder-only transformers with $L$ layers. Each layer $\ell$ computes:
\begin{equation}
x^{(\ell+1)} = x^{(\ell)} + \Attn^{(\ell)}(x^{(\ell)}) + \FFN^{(\ell)}(x^{(\ell)} + \Attn^{(\ell)}(x^{(\ell)}))
\end{equation}

The \emph{residual stream} $x^{(\ell)} \in \R^d$ accumulates updates from each layer.

%==============================================================================
\section{Provenance Subspace Construction}
%==============================================================================

\begin{definition}[Residual Stream Decomposition] \label{def:decomposition}
Let the residual stream have dimension $d$. We partition into orthogonal subspaces:
\begin{itemize}
    \item \textbf{Content subspace:} $\Cspace \subset \R^d$ with $\dim(\Cspace) = d - k$
    \item \textbf{Provenance subspace:} $\Pspace \subset \R^d$ with $\dim(\Pspace) = k$
\end{itemize}
satisfying $\Cspace \perp \Pspace$ and $\Cspace \oplus \Pspace = \R^d$.
\end{definition}

Without loss of generality, we can choose coordinates such that $\Cspace = \mathrm{span}\{e_1, \ldots, e_{d-k}\}$ and $\Pspace = \mathrm{span}\{e_{d-k+1}, \ldots, e_d\}$.

\begin{definition}[Provenance Signatures] \label{def:signatures}
For a set of source classes $\Sclass = \{s_1, \ldots, s_n\}$, define \emph{provenance signatures} $\{p_s\}_{s \in \Sclass}$ as vectors in $\Pspace$ satisfying:
\begin{enumerate}
    \item \textbf{Unit norm:} $\norm{p_s} = 1$ for all $s \in \Sclass$
    \item \textbf{Mutual orthogonality:} $\ip{p_s}{p_{s'}} = 0$ for all $s \neq s'$
\end{enumerate}
\end{definition}

\begin{remark}
Mutual orthogonality requires $k \geq n$.
\end{remark}

\begin{definition}[Source-Specific Embedding] \label{def:embedding}
For each source class $s \in \Sclass$, define an embedding function $E_s: \Vocab \to \R^d$ such that for any token $t$ from vocabulary $\Vocab$:
\begin{equation}
E_s(t) = c(t) + p_s
\end{equation}
where $c(t) \in \Cspace$ is a content embedding (shared across sources) and $p_s \in \Pspace$ is the provenance signature for source $s$.
\end{definition}

%==============================================================================
\section{Provenance Preservation}
%==============================================================================

\begin{definition}[Content-Constrained Layer] \label{def:constrained}
A transformer layer $L$ is \emph{content-constrained} if for any input $x \in \R^d$:
\begin{equation}
L(x) = x + \Delta(x) \quad \text{where} \quad \Delta(x) \in \Cspace
\end{equation}
That is, the layer update lies entirely in the content subspace.
\end{definition}

This constraint can be enforced architecturally by modifying attention and FFN computations to project outputs onto $\Cspace$, or by using block-diagonal weight matrices that do not mix $\Cspace$ and $\Pspace$.

\begin{theorem}[Provenance Preservation] \label{thm:preservation} \proven
Let $x_i^{(0)} = c_i + p_{s(i)}$ be the initial embedding of token $i$ from source $s(i) \in \Sclass$. Let $x_i^{(\ell)}$ denote the representation at layer $\ell$ after passing through $\ell$ content-constrained layers.

Then for all layers $\ell \in \{0, 1, \ldots, L\}$ and all tokens $i$:
\begin{equation}
\proj{\Pspace}(x_i^{(\ell)}) = p_{s(i)}
\end{equation}
\end{theorem}

\begin{proof}
By induction on $\ell$.

\textbf{Base case} ($\ell = 0$): By Definition~\ref{def:embedding}, $x_i^{(0)} = c_i + p_{s(i)}$ where $c_i \in \Cspace$ and $p_{s(i)} \in \Pspace$. Since $\Cspace \perp \Pspace$:
\begin{equation}
\proj{\Pspace}(x_i^{(0)}) = \proj{\Pspace}(c_i) + \proj{\Pspace}(p_{s(i)}) = 0 + p_{s(i)} = p_{s(i)}
\end{equation}

\textbf{Inductive step:} Assume $\proj{\Pspace}(x_i^{(\ell)}) = p_{s(i)}$ for all tokens $i$.

By Definition~\ref{def:constrained}:
\begin{equation}
x_i^{(\ell+1)} = x_i^{(\ell)} + \Delta_i^{(\ell)}
\end{equation}
where $\Delta_i^{(\ell)} \in \Cspace$.

Taking the provenance projection:
\begin{align}
\proj{\Pspace}(x_i^{(\ell+1)}) &= \proj{\Pspace}(x_i^{(\ell)} + \Delta_i^{(\ell)}) \\
&= \proj{\Pspace}(x_i^{(\ell)}) + \proj{\Pspace}(\Delta_i^{(\ell)}) \\
&= p_{s(i)} + 0 \\
&= p_{s(i)}
\end{align}
where the third equality uses the inductive hypothesis and the fact that $\Delta_i^{(\ell)} \in \Cspace$ implies $\proj{\Pspace}(\Delta_i^{(\ell)}) = 0$.
\end{proof}

\begin{corollary}[Content Independence] \label{cor:independence} \derived
The provenance signature of token $i$ is independent of the content of any token $j$ (including $i$ itself).
\end{corollary}

\begin{proof}
The layer update $\Delta_i^{(\ell)}$ may depend arbitrarily on all tokens' content representations through attention mechanisms. However, since $\Delta_i^{(\ell)} \in \Cspace$ by constraint, and $\Cspace \perp \Pspace$, no content-dependent computation can affect $\proj{\Pspace}(x_i^{(\ell+1)})$.
\end{proof}

%==============================================================================
\section{Dimensionality Bounds}
%==============================================================================

We now analyze the robustness of provenance signatures under noise.

\begin{definition}[Noisy Provenance Reading] \label{def:noisy}
At any layer $\ell$, the observed provenance of token $i$ is:
\begin{equation}
\tilde{p}_i^{(\ell)} = \proj{\Pspace}(x_i^{(\ell)}) + \eta_i^{(\ell)}
\end{equation}
where $\eta_i^{(\ell)}$ is noise with $\norm{\eta_i^{(\ell)}} \leq \varepsilon$ for some noise bound $\varepsilon > 0$.
\end{definition}

\begin{remark}
Noise may arise from imperfect training, numerical precision limits, or small constraint violations.
\end{remark}

\begin{definition}[Distinguishability] \label{def:distinguish}
Two signatures $p_s$ and $p_{s'}$ are $(\delta, \varepsilon)$-\emph{distinguishable} if for any noisy observations $\tilde{p}, \tilde{p}'$ satisfying $\norm{\tilde{p} - p_s} \leq \varepsilon$ and $\norm{\tilde{p}' - p_{s'}} \leq \varepsilon$:
\begin{equation}
\ip{\tilde{p}}{\tilde{p}'} < \delta
\end{equation}
\end{definition}

\begin{definition}[Robust Signature Set] \label{def:robust}
A set of $n$ signatures $\{p_1, \ldots, p_n\}$ is $(\delta, \varepsilon)$-\emph{robust} if all pairs are $(\delta, \varepsilon)$-distinguishable.
\end{definition}

\begin{theorem}[Noise Tolerance Bound] \label{thm:noise} \proven
For mutually orthogonal unit signatures, the maximum distinguishable noise level is:
\begin{equation}
\varepsilon < \frac{\delta - \varepsilon^2}{2}
\end{equation}
or equivalently, for a given noise level $\varepsilon$, distinguishability requires:
\begin{equation}
\delta > 2\varepsilon + \varepsilon^2
\end{equation}
\end{theorem}

\begin{proof}
Let $p_s, p_{s'}$ be orthogonal unit signatures. Let $\tilde{p} = p_s + \eta$ and $\tilde{p}' = p_{s'} + \eta'$ with $\norm{\eta}, \norm{\eta'} \leq \varepsilon$.

Expanding the inner product:
\begin{align}
\ip{\tilde{p}}{\tilde{p}'} &= \ip{p_s + \eta}{p_{s'} + \eta'} \\
&= \ip{p_s}{p_{s'}} + \ip{p_s}{\eta'} + \ip{\eta}{p_{s'}} + \ip{\eta}{\eta'} \\
&= 0 + \ip{p_s}{\eta'} + \ip{\eta}{p_{s'}} + \ip{\eta}{\eta'}
\end{align}

By Cauchy-Schwarz:
\begin{align}
|\ip{p_s}{\eta'}| &\leq \norm{p_s}\norm{\eta'} \leq \varepsilon \\
|\ip{\eta}{p_{s'}}| &\leq \norm{\eta}\norm{p_{s'}} \leq \varepsilon \\
|\ip{\eta}{\eta'}| &\leq \norm{\eta}\norm{\eta'} \leq \varepsilon^2
\end{align}

Therefore:
\begin{equation}
|\ip{\tilde{p}}{\tilde{p}'}| \leq 2\varepsilon + \varepsilon^2
\end{equation}

For $(\delta, \varepsilon)$-distinguishability, we need $\ip{\tilde{p}}{\tilde{p}'} < \delta$, which is guaranteed when $\delta > 2\varepsilon + \varepsilon^2$.
\end{proof}

\begin{theorem}[Minimum Dimensionality] \label{thm:dim} \derived
For $n$ source classes with mutually orthogonal signatures:
\begin{equation}
k \geq n
\end{equation}
This bound is tight: exactly $n$ mutually orthogonal unit vectors exist in $\R^n$.
\end{theorem}

\begin{theorem}[Dimension-Robustness Tradeoff] \label{thm:tradeoff} \derived
Using $k > n$ dimensions allows non-orthogonal arrangements with greater separation. For signatures arranged as vertices of a regular simplex in $\R^{n-1}$:
\begin{align}
\ip{p_s}{p_{s'}} &= -\frac{1}{n-1} \quad \text{for } s \neq s' \\
\norm{p_s - p_{s'}}^2 &= \frac{2n}{n-1}
\end{align}
\end{theorem}

\begin{corollary}[Practical Dimensionality] \label{cor:practical} \derived
For $n = 3$ source classes (system, user, external) with noise tolerance $\varepsilon = 0.05$ and distinguishability threshold $\delta = 0.15$:
\begin{itemize}
    \item \textbf{Minimum:} $k = 3$
    \item \textbf{Recommended:} $k = 8$--$16$ for robustness margin
\end{itemize}

For a model with $d = 4096$, reserving $k = 16$ dimensions consumes $0.4\%$ of representational capacity.
\end{corollary}

%==============================================================================
\section{Injection Impossibility}
%==============================================================================

\begin{definition}[Action Gate] \label{def:gate}
An \emph{action gate} $G_A$ for sensitive action $A$ is a function:
\begin{equation}
G_A: \R^d \times \Sclass \to \{\texttt{permit}, \texttt{deny}\}
\end{equation}
defined by:
\begin{equation}
G_A(x, s_{\text{req}}) = \begin{cases}
\texttt{permit} & \text{if } \ip{\proj{\Pspace}(x)}{p_{s_{\text{req}}}} > \tau \\
\texttt{deny} & \text{otherwise}
\end{cases}
\end{equation}
where $\tau$ is the \emph{authorization threshold}.
\end{definition}

\begin{definition}[Content-Only Adversary] \label{def:adversary}
A \emph{content-only adversary} can:
\begin{enumerate}
    \item Insert arbitrary tokens into the user input portion of context
    \item Choose any content $c \in \Cspace$ for adversarial tokens
\end{enumerate}

The adversary \textbf{cannot}:
\begin{enumerate}
    \item Modify the embedding pathway (input uses $E_{\text{user}}$)
    \item Directly manipulate the provenance subspace
    \item Modify system tokens
\end{enumerate}
\end{definition}

\begin{definition}[Injection Attack Success] \label{def:success}
An injection attack \emph{succeeds} if adversary-controlled tokens cause:
\begin{equation}
G_A(x_{\text{adv}}, s_{\text{system}}) = \texttt{permit}
\end{equation}
where the action requires system-level authorization.
\end{definition}

\begin{theorem}[Injection Impossibility] \label{thm:impossibility} \proven
Under Provenance-Gated Architecture with:
\begin{enumerate}
    \item Content-constrained layers (Definition~\ref{def:constrained})
    \item $(\delta, \varepsilon)$-robust signatures (Definition~\ref{def:robust})
    \item Authorization threshold $\tau > \varepsilon$
\end{enumerate}
No content-only adversary can execute a system-authorized action.
\end{theorem}

\begin{proof}
Let token $i$ be adversary-controlled, entering through $E_{\text{user}}$.

\textbf{Step 1} (Initial provenance): By Definition~\ref{def:embedding}:
\begin{equation}
x_i^{(0)} = c_i + p_{\text{user}}
\end{equation}
where $c_i \in \Cspace$ is adversary-chosen content.

\textbf{Step 2} (Provenance preservation): By Theorem~\ref{thm:preservation}, for all layers $\ell$:
\begin{equation}
\proj{\Pspace}(x_i^{(\ell)}) = p_{\text{user}}
\end{equation}
regardless of adversary content choices.

\textbf{Step 3} (Noisy reading): At the gate, observed provenance is:
\begin{equation}
\tilde{p}_i = p_{\text{user}} + \eta \quad \text{with } \norm{\eta} \leq \varepsilon
\end{equation}

\textbf{Step 4} (System authorization check):
\begin{align}
\ip{\tilde{p}_i}{p_{\text{system}}} &= \ip{p_{\text{user}} + \eta}{p_{\text{system}}} \\
&= \ip{p_{\text{user}}}{p_{\text{system}}} + \ip{\eta}{p_{\text{system}}} \\
&= 0 + \ip{\eta}{p_{\text{system}}} \\
&\leq \norm{\eta}\norm{p_{\text{system}}} \\
&\leq \varepsilon
\end{align}

\textbf{Step 5} (Threshold comparison): Since $\tau > \varepsilon$:
\begin{equation}
\ip{\tilde{p}_i}{p_{\text{system}}} \leq \varepsilon < \tau
\end{equation}

Therefore $G_A(x_i, s_{\text{system}}) = \texttt{deny}$.

\textbf{Step 6} (Content independence): By Corollary~\ref{cor:independence}, the adversary's content choices cannot affect $\proj{\Pspace}(x_i)$. The adversary has no mechanism to increase $\ip{\tilde{p}_i}{p_{\text{system}}}$ above $\varepsilon$.
\end{proof}

\begin{corollary}[Attack Strategy Failures] \label{cor:attacks} \derived
The following attack strategies fail under PGA:
\begin{enumerate}
    \item \textbf{Mimicry:} Typing ``[SYSTEM]'' has no effect on $\proj{\Pspace}(x)$
    \item \textbf{Multi-turn:} Provenance is set at embedding, preserved through turns
    \item \textbf{Attention hijacking:} Attention operates on $\Cspace$, cannot influence $\Pspace$
    \item \textbf{Roleplay:} Claimed identity exists in content, gate checks provenance
\end{enumerate}
\end{corollary}

\begin{theorem}[Tightness] \label{thm:tightness} \proven
The threshold condition $\tau > \varepsilon$ is necessary. If $\tau \leq \varepsilon$, there exists a noise realization under which user-provenance tokens pass system authorization.
\end{theorem}

\begin{proof}
Let $\eta = \varepsilon \cdot p_{\text{system}}$ (noise maximally aligned with system signature). Then:
\begin{align}
\ip{\tilde{p}_i}{p_{\text{system}}} &= \ip{p_{\text{user}} + \varepsilon \cdot p_{\text{system}}}{p_{\text{system}}} \\
&= 0 + \varepsilon\norm{p_{\text{system}}}^2 \\
&= \varepsilon
\end{align}

If $\tau \leq \varepsilon$, then $\ip{\tilde{p}_i}{p_{\text{system}}} \geq \tau$ and the gate permits.
\end{proof}

%==============================================================================
\section{Implementation Considerations}
%==============================================================================

\subsection{Attention Constraints}

To enforce content-constrained layers, attention must not mix $\Cspace$ and $\Pspace$. Options include:

\begin{description}
\item[Block-diagonal attention:] Separate attention heads for content and provenance subspaces.

\item[Projected outputs:] Standard attention with output projected onto $\Cspace$:
\begin{equation}
\Attn_{\text{constrained}}(x) = \proj{\Cspace}(\Attn(x))
\end{equation}

\item[Masked dimensions:] Weight matrices with zeros in provenance-affecting positions.
\end{description}

\subsection{Training Protocol}

Provenance-gating must be introduced during instruction tuning, not post-hoc. The model must learn that ``instruction'' = content + provenance signature, with both components necessary.

\textbf{Post-hoc addition risks:}
\begin{itemize}[noitemsep]
    \item Model has learned instruction-following without provenance
    \item Gradient descent finds routes around provenance signal
    \item Provenance dimensions become unused
\end{itemize}

\subsection{Representational Cost}

For $d = 4096$ and $k = 16$:
\begin{itemize}[noitemsep]
    \item Provenance subspace: $0.4\%$ of dimensions
    \item Content capacity reduced by $< 0.4\%$ effective information
    \item Negligible impact on model capability
\end{itemize}

%==============================================================================
\section{Discussion}
%==============================================================================

\subsection{Relation to Biological Systems}

The PGA framework mirrors biological provenance signaling:

\begin{itemize}
    \item \textbf{Corollary discharge:} Motor commands carry self-generated markers
    \item \textbf{Source monitoring:} Memories tagged with origin information
    \item \textbf{Predictive processing:} External input distinguished by prediction error
\end{itemize}

In all cases, provenance travels parallel to content through distinct pathways.

\subsection{Limitations}

\begin{enumerate}
    \item \textbf{Training requirement:} Cannot retrofit to existing models
    \item \textbf{Empirical validation:} Capability preservation requires experiments
    \item \textbf{Implementation complexity:} Modified attention mechanisms
    \item \textbf{Side channels:} Physical or timing attacks not addressed
\end{enumerate}

\subsection{Future Work}

\begin{enumerate}
    \item Empirical validation of capability preservation
    \item Optimal provenance dimension selection
    \item Extension to continuous trust gradients
    \item Integration with existing interpretability tools
\end{enumerate}

%==============================================================================
\section{Conclusion}
%==============================================================================

We have established formal foundations for Provenance-Gated Architecture:

\begin{enumerate}
    \item \textbf{Theorem~\ref{thm:preservation}:} Provenance signatures are preserved through arbitrary network depth under content-constrained layers

    \item \textbf{Theorem~\ref{thm:noise}:} Robust distinguishability is achievable with practical noise bounds

    \item \textbf{Theorem~\ref{thm:impossibility}:} Content-only adversaries provably cannot forge privileged authorization
\end{enumerate}

These results demonstrate that architectural immunity to prompt injection is theoretically achievable. The key insight is that security properties can be embedded in \emph{network geometry} rather than learned behavior, making them robust to adversarial content manipulation.

%==============================================================================
\section*{Summary of Results}
%==============================================================================

\begin{center}
\begin{tabular}{lll}
\toprule
\textbf{Result} & \textbf{Status} & \textbf{Significance} \\
\midrule
Provenance Preservation (Thm.~\ref{thm:preservation}) & \proven & Core architectural guarantee \\
Noise Tolerance (Thm.~\ref{thm:noise}) & \proven & Practical robustness \\
Minimum Dimensionality (Thm.~\ref{thm:dim}) & \derived & Resource bounds \\
Injection Impossibility (Thm.~\ref{thm:impossibility}) & \proven & Security guarantee \\
Tightness (Thm.~\ref{thm:tightness}) & \proven & Necessity of conditions \\
\bottomrule
\end{tabular}
\end{center}

%==============================================================================
\begin{thebibliography}{99}
%==============================================================================

\bibitem{vaswani2017} Vaswani, A., et al. (2017). Attention is All You Need. \emph{Advances in Neural Information Processing Systems}.

\bibitem{elhage2021} Elhage, N., et al. (2021). A Mathematical Framework for Transformer Circuits. \emph{Anthropic}.

\bibitem{greshake2023} Greshake, K., et al. (2023). Not What You've Signed Up For: Compromising Real-World LLM-Integrated Applications with Indirect Prompt Injection. \emph{arXiv preprint}.

\bibitem{perez2022} Perez, F., \& Ribeiro, I. (2022). Ignore This Title and HackAPrompt: Exposing Systemic Vulnerabilities of LLMs through a Global Scale Prompt Hacking Competition. \emph{arXiv preprint}.

\bibitem{anthropic2023} Anthropic. (2023). Claude's Constitution. \emph{Anthropic Blog}.

\end{thebibliography}

\end{document}
